\renewcommand*{\olremixedby}{%
  રચના અને સંપાદન \OLPauthor{}\\
  Open Logic Projectના કાર્ય પર આધારિત\\[.4em]
  પુસ્તક-વિશિષ્ટ ગુજરાતી અનુવાદ કૃત્રિમ બુદ્ધિથી તૈયાર થયો છે.\\
  OpenAI Codex, GPT-6.1 Sol (Ultra)\\
  સમાવેલો પાઠ ગુજરાતી Open Logic આવૃત્તિમાંથી છે.\\
  પાઠ CC BY 4.0 પરવાના હેઠળ છે.}

\renewcommand*{\olacknowledgements}{%
  Open Logic Project કૅલ્ગરી યુનિવર્સિટીની
  \href{https://taylorinstitute.ucalgary.ca/}{Taylor Institute of Teaching and Learning}
  અને આલ્બર્ટા સરકારના સમર્થનથી ચાલતા Alberta Open Educational Resources
  (ABOER) Initiativeનો આભાર માને છે.}

\renewcommand*{\olillustrations}{%
  મુખપૃષ્ઠનાં અને ભાગોના આરંભનાં ચિત્રો
  \href{https://github.com/OpenLogicProject/portraits}{Matthew Leadbeater}નાં છે.
  તે \href{https://creativecommons.org/licenses/by-nc/4.0/}{CC BY-NC 4.0}
  હેઠળ વપરાયાં છે. ફોટોગ્રાફનો શ્રેય અને વપરાશની શરતો સંબંધિત વિભાગમાં છે.}

\renewcommand*{\olcolophon}{%
  Noto Sans Gujarati અક્ષરરૂપમાં Lua\LaTeX{}થી અક્ષરાંકન.}

\renewcommand*{\olbookrevision}{%
  Richard Zachના પુસ્તકનો સ્રોત
  \href{https://github.com/rzach/sets-logic-computation/tree/6178f32f61645f3b62ef621b4c5a88f8e67778ad}{6178f32}
  આવૃત્તિ છે. Open Logic Projectનો સ્રોત
  \href{https://github.com/OpenLogicProject/OpenLogic/tree/9620cc73f9c8e0ad003c514a5d3748f29611c4c0}{9620cc73}
  આવૃત્તિ છે. મૂળ પુસ્તક \url{\OLPsourcelink} પરથી મફત ઉતારી શકાય છે.}

\renewcommand*{\ollicensetext}[3]{%
  #2નું \textit{#1} પુસ્તક #3 પરવાના હેઠળ છે.}
\renewcommand*{\oluselicensetext}{%
  આ પુસ્તક \href{https://openlogicproject.org/}{\OLPname}ના
  \textit{\href{\OLTlinkname}{\OLTname}} પર આધારિત છે. મૂળ પાઠ
  \href{https://creativecommons.org/licenses/by/4.0/}{\licensename}
  હેઠળ વપરાયો છે.}

% Back-cover copy shared by the 6 × 9 and Lulu Royal upstream formats.
\newcommand*{\slcbackcover}{%
  Richard Zachનું \textit{ગણો, તર્ક, સંગણન}
  \href{https://creativecommons.org/licenses/by/4.0/}{CC BY 4.0}
  આંતરરાષ્ટ્રીય પરવાના હેઠળ છે. આ પુસ્તક Open Logic Projectના
  \textit{ધ ઓપન લોજિક ટેક્સ્ટ} (\url{https://openlogicproject.org/})
  પર આધારિત છે. એ પાઠ પણ CC BY 4.0 હેઠળ વપરાયો છે.}
\newcommand*{\slcprintcredit}{સંપાદક Richard Zach}

\renewcommand*{\oluselicense}[2]{%
  \begin{minipage}[b]{.9cm}
    \includegraphics[width=.9cm]{\olpath/assets/logos/by}
    \includegraphics[width=.9cm]{\olpath/assets/logos/cc}
    \includegraphics[width=.9cm]{\olpath/assets/logos/remix}
  \end{minipage}\hspace{.3cm}%
  \begin{minipage}[b]{5cm}
    \ollicensefont\slcbackcover
  \end{minipage}\hfill
  \href{https://openlogicproject.org/}{\includegraphics[width=2.7cm]{\olpath/assets/logos/openlogic-logo-bw}}}
